\CatchFileBetweenTags{\AlphaInvVal}{constants.tex}{AlphaInvVal}
\CatchFileBetweenTags{\AlphaInvMAR}{constants.tex}{AlphaInvMAR}
\CatchFileBetweenTags{\MeMeVPrint}{constants.tex}{MeMeVPrint}

% CAP
\CatchFileBetweenTags{\HiggsVEVStepOneVal}{constants.tex}{HiggsVEVStepOneVal}
\CatchFileBetweenTags{\HiggsVEVStepTwoVal}{constants.tex}{HiggsVEVStepTwoVal}
\CatchFileBetweenTags{\HiggsVEVVal}{constants.tex}{HiggsVEVVal}
\CatchFileBetweenTags{\HiggsVEVExperimentalValue}{constants.tex}{HiggsVEVExperimentalValue}
\CatchFileBetweenTags{\HiggsVEVSigma}{constants.tex}{HiggsVEVSigma}
\CatchFileBetweenTags{\HiggsImpedanceVal}{constants.tex}{HiggsImpedanceVal}
\CatchFileBetweenTags{\WeakApertureProjVal}{constants.tex}{WeakApertureProjVal}

\CatchFileBetweenTags{\FermiConstVal}{constants.tex}{FermiConstVal}
\CatchFileBetweenTags{\FermiConstExperimentalValue}{constants.tex}{FermiConstExperimentalValue}
\CatchFileBetweenTags{\FermiConstSigma}{constants.tex}{FermiConstSigma}

% PRO
\CatchFileBetweenTags{\HiggsLambdaVal}{constants.tex}{HiggsLambdaVal}
\CatchFileBetweenTags{\HiggsLambdaStepOneVal}{constants.tex}{HiggsLambdaStepOneVal}
\CatchFileBetweenTags{\HiggsLambdaExperimentalValue}{constants.tex}{HiggsLambdaExperimentalValue}
\CatchFileBetweenTags{\HiggsLambdaSigma}{constants.tex}{HiggsLambdaSigma}

% TOL
\CatchFileBetweenTags{\HiggsMassVal}{constants.tex}{HiggsMassVal}
\CatchFileBetweenTags{\HiggsMassExperimentalValue}{constants.tex}{HiggsMassExperimentalValue}
\CatchFileBetweenTags{\HiggsMassSigma}{constants.tex}{HiggsMassSigma}

% MAR
\CatchFileBetweenTags{\ElectronYukawaVal}{constants.tex}{ElectronYukawaVal}
\CatchFileBetweenTags{\ElectronYukawaStepOneVal}{constants.tex}{ElectronYukawaStepOneVal}
\CatchFileBetweenTags{\ElectronYukawaExperimentalValue}{constants.tex}{ElectronYukawaExperimentalValue}
\CatchFileBetweenTags{\ElectronYukawaSigma}{constants.tex}{ElectronYukawaSigma}

\section{System 5: The Resonant Interface (The Electroweak Scale)}
\label{sec:resonant_interface}

The vacuum possesses a strictly finite channel capacity ($\nu=16$) operating at a fundamental temporal resonance ($\Delta=43$). In physical terms, this finite bandwidth imposes a strict ceiling on the local energy density of the lattice ($E = \hbar\omega$). The introduction of complex Topological Defects into the bulk saturates this core capacity, threatening the vacuum with Causal Aliasing (runaway energy divergence). 

To persist, the system must resolve this complexity. While System 4 defined the \textit{fractional ratio} of bandwidth allocated to this sector (e.g., $\sin^2\theta_W$), System 5 derives the \textit{absolute magnitude} of that capacity. This absolute ceiling manifests as the Electroweak Scale.

Under \textbf{Rule 3: Impedance Matching}, the vacuum must translate information between the deep geometric bulk (the core lattice) and the localized topological boundary (the particle surface) without dissipative reflection loss. To achieve this, the active energy bandwidth of the Weak Interaction must be perfectly matched to the bulk lattice resonance. Here, we derive the exact phenomenological constants of this interface as the partition coefficients of this geometric matching.

Consistent with \textit{Informational Energetics: Entropic Action} \cite{meyer_ieaction_2026}, we identify the Higgs sector as the mathematical bounding of this state-space capacity. The internal bandwidth of this resonant interface is rigidly partitioned across the six pillars, defining the following physical parameters:

\begin{itemize}
    \item \textbf{Capacity ($v$):} The VEV; the mass baseline allocated to the resonant interface, against which the vacuum potential is normalized.
    \item \textbf{Map ($Y$):} Scalar Hypercharge; the unitary address. It corresponds to the inverse of the lattice Unitary Ground State ($\Delta^0=1$).
    \item \textbf{Protocol ($\lambda$):} Self-Coupling; the self-interaction bandwidth allocation.
    \item \textbf{Governor ($\lambda|\phi|^4$):} The Quartic Potential; The geometric capacity wall enforcing the topological boundary ($\chi=2$).
    \item \textbf{Toll ($\mu^2$):} Instability Factor; the fractional entropic cost of maintaining the broken symmetry state.
    \item \textbf{Margin ($y_e$):} The Electron Yukawa; the minimal resolvable noise floor.
\end{itemize}

In the Standard Model, the seemingly arbitrary values of this sector constitute the ``Hierarchy Problem.'', but they are recognized here as the unique set of partition coefficients required to maintain vacuum stability against causal aliasing.

\subsection{Capacity: The Higgs VEV (\texorpdfstring{$v$}{v})}
Under \textbf{Rule 1: Capacity Partitioning}, the VEV represents the net informational capacity allocated to the Weak Interaction. We derive this by calculating the total available states over the causal boundary and subtracting the structural overhead.

\subsubsection{The Bare Geometric Floor}
To derive the absolute capacity of the Electroweak interface, we must first define the total informational envelope required to sustain a stable topological defect. Under \textbf{Rule 1: Capacity Partitioning}, the system must allocate enough bandwidth to account for every state transition within the defect's causal influence.

The temporal period $\Delta=43$ defines the \textit{causal horizon} of a lattice interaction (the distance information can travel before the system resets). Because information propagates isotopically at the speed of light ($c=1$), this resonance defines a 2-dimensional causal area $\Delta \times \Delta = \Delta^2$. 

To distinguish a specific particle (a defect) from the vacuum, this area must be bounded by the topological invariant $\chi=2$ (the sphere). The total capacity envelope required to resolve the defect's identity against the background is therefore the product of its topological signature and its resonant area: $\chi \Delta^2$. This represents the total number of ``state-slots'' available to the interface over one cycle. 

From this total envelope, the substrate must reserve the total information the vacuum must process per fundamental cycle: ($L_{substrate} = \nu + D \cdot \Delta = 188$\cite{meyer_ieaction_2026}). Multiplying by the baseline vacuum impedance ($\alpha^{-1}$) we can get the transmission impedance.

This transmission impedance is a strictly dimensionless ratio which can be mapped into a physical energy scale by anchoring it to the geometric resolution floor of the lattice: the minimum viable bit of mass ($m_e$, derived structurally below in Section \ref{subsec:margin_electron}). This establishes the true physical energy baseline of the field \textit{before} it is subjected to the macroscopic topological dilution and thermal noise corrections ($\delta_{topo}, \delta_{noise}$).

\begin{equation}
\begin{split}
v_{geo} &= (\chi \Delta^2 - L_{substrate}) \cdot \alpha^{-1} \cdot m_e \\
&= (2 \cdot 43^2 - 188) \cdot \AlphaInvVal \cdot \MeMeVPrint \text{ MeV} \\
&\approx \HiggsVEVStepOneVal
\end{split}
\end{equation}

\subsubsection{The Geometric Corrections}
The $v_{geo}$ baseline represents an idealized count of state-slots. To find the \textit{physical} energy density, we must account for the reality of the projection: how this energy is distributed in continuous space and how it is resolved against the substrate noise floor. These effects manifest as exact leading-order linear expansions ($1 \pm x$), representing the first-order response of the vacuum to its own geometric constraints.

\paragraph{Topological Dilution (Spatial Correction)}
In the continuum limit, the localized charge of a defect is not a discrete point but a field distributed over a spherical gauge geometry ($4\pi$). This ``dilutes'' the available energy density because the topological boundary ($\chi=2$) effectively expands the volume the coupling must cover. We define the \textbf{Effective Coupling Volume} ($D_{eff}$) as the manifold rank ($D=4$) augmented by this boundary-to-volume ratio:
\begin{equation}
D_{eff} = D + \frac{\chi}{4\pi} \approx 4.15915
\end{equation}
\textit{(Scope Defense: This continuous $4\pi$ correction represents the macroscopic solid angle of the $\chi=2$ boundary; it does not alter the integer dimensionality of the substrate).} This yields the dilution factor:
\begin{equation}
\delta_{topo} = \left( 1 + \frac{\alpha}{D_{eff}} \right)
\end{equation}

\paragraph{Thermodynamic Noise (Temporal Correction)}
We must also subtract the energy ``lost'' to the substrate's resolution floor. Under \textbf{Rule 2: Reversible Encoding}, the \textbf{Margin Impedance} ($Z_{MAR} \approx \AlphaInvMAR$) defines the minimum resolvable energy unit. Because the vacuum supports exactly 3 fermion generations (the available residual degrees of freedom: $\sigma - \chi = 3$), this noise is partitioned across these three available ``signal slots.'' To maintain global solvency, the physical VEV must be reduced by this noise amplitude:
\begin{equation}
\delta_{noise} = \left( 1 - \frac{Z_{MAR}}{\sigma - \chi} \right)
\end{equation}

\subsubsection{Synthesis: The Physical VEV}
The physical Vacuum Expectation Value is the geometric floor modulated by the simultaneous application of these topological and thermodynamic factors.
\begin{equation}
v_{phys} = v_{geo} \cdot \delta_{topo} \cdot \delta_{noise}
\end{equation}

Substituting the derived values:
\begin{equation}
\begin{split}
v_{phys} &= (\HiggsVEVStepOneVal) \cdot \\ 
&\left( 1 + \frac{\alpha}{4.159} \right) \cdot \left( 1 - \frac{\AlphaInvMAR}{3} \right) \\
& \approx \mathbf{\HiggsVEVVal}
\end{split}
\end{equation}

\vspace{0.5em}
\begin{tabular}{ll}
\textbf{Prediction:} & $\HiggsVEVVal$ \\
\textbf{Experiment:} & \HiggsVEVExperimentalValue \\
\textbf{Accuracy:}   & $\HiggsVEVSigma \sigma$ \\
\end{tabular}
\vspace{0.5em}

\subsubsection{Validation: The Fermi Constant (\texorpdfstring{$G_F$}{GF})}
To validate the derived energy capacity of the interface ($v_{phys}$), we calculate the \textbf{Fermi Constant} ($G_F$), which serves as the experimental measure of the Weak interaction's effective strength. In standard QFT, these are related by $G_F = 1/(\sqrt{2}v^2)$. While the $\sqrt{2}$ is treated as a historical artifact in standard physics, it emerges here as a necessary geometric projection.

\paragraph{The Geometric Origin of $\sqrt{2}$}
In Enrico Fermi’s original theory of beta decay, he assumed parity was conserved. When the discovery of maximal parity violation later proved that only left-handed neutrinos participate in the Weak interaction, the $\sqrt{2}$ was introduced as a manual scale correction to align the theory with observed data.

In $E_8$-Persistence theory, for the $v_{phys}$ to mediate a force it must be projected onto the 1D chiral trajectory (the direction of travel) of a fermion. The geometric projection of a 2-dimensional isotropic amplitude onto a 1-dimensional axis is defined by the root-mean-square (RMS) projection, which yields exactly $1/\sqrt{2}$.

\paragraph{The Result}
The Fermi constant is therefore the inverse squared capacity of the interface, attenuated by this geometric projection factor:
\begin{equation}
    G_F = \frac{1}{\sqrt{2}(v_{phys})^2} \approx \mathbf{\FermiConstVal}
\end{equation}


\vspace{0.5em}
\begin{tabular}{ll}
\textbf{Prediction:} & $\FermiConstVal$ \\
\textbf{Experiment:} & \FermiConstExperimentalValue \\
\textbf{Accuracy:}   & $\FermiConstSigma \sigma$ \\
\end{tabular}
\vspace{0.5em}

\subsection{Map: Scalar Hypercharge (\texorpdfstring{$Y=1$}{Y})} 
To coordinate information across the interface, the regulator requires a discrete \textbf{Map} address. As established in \cite{meyer_ieaction_2026}, the ``origin'' coordinate of the lattice is the Unitary Ground State ($\Delta^0 = 1$). To maintain a perfect resonance with the vacuum, the regulator's address must correspond to this fundamental unitary charge.

Combined with the topological boundary ($\chi=2$), this unitary address geometrically necessitates the $SU(2)$ doublet structure found in standard physics. This allows the regulator to mediate between the chiral matter sector and the gauge field bulk. (\textit{Note: A complete derivation of how all Standard Model fermion hypercharges emerge as exact geometric ratios of these lattice invariants is provided in Appendix \ref{app:OriginOfHypercharge}}).



\subsection{Protocol: Self-Coupling (\texorpdfstring{$\lambda$}{lambda})}
In standard QFT, the Higgs self-coupling ($\lambda$) is an arbitrary parameter that determines the "stiffness" of the vacuum potential. In the $E_8$-Persistence theory, this is the \textbf{Protocol} for local self-regulation. Under \textbf{Rule 1: Capacity Partitioning}, $\lambda$ is derived as the exact fraction of the defect's bandwidth dedicated to maintaining its own internal stability.

\subsubsection{The Intrinsic Structural Load (\texorpdfstring{$L_{intrinsic}$}{Lintrinsic})}
The ``cost'' of the defect is a fraction of by its total informational complexity. As established in \cite{meyer_informational_2026}, a defect's total informational complexity is its \textbf{Intrinsic Load} $(L_{intrinsic} = 23)$ is the sum of its orthogonal geometric components: its chiral capacity ($\nu=16$), its interaction degrees of freedom ($\sigma=5$), and its topological identity lock ($\chi=2$).

\subsubsection{The Regulatory Bandwidth}
The ``budget'' available for self-regulation is the capacity remaining in the internal interaction bulk ($\sigma - \chi = 3$) after the identity lock is paid. However, the system must pay a \textbf{Resonant Cost} ($1/\Delta$) to maintain temporal synchronization with the lattice. Without reserving this fractional bandwidth to ``keep time,'' the scalar field would drift out of phase with the fundamental update cycle ($\Delta=43$), causing the subsystem to decohere and decay.

The net bandwidth available for self-regulation is therefore:
\begin{equation}
\begin{split}    
\text{Net} &= (\sigma - \chi) - \frac{1}{\Delta} \\
&= 3 - \frac{1}{43} \approx \HiggsLambdaStepOneVal
\end{split}
\end{equation}

\subsubsection{The Resulting Coupling}
The coupling $\lambda$ is the ratio of this regulatory budget to the total structural load. It represents the ``stiffness'' required to anchor the defect against the vacuum's entropic current:
\begin{equation}
\lambda = \frac{\text{Net}}{L_{intrinsic}} = \frac{2.9767}{23} \approx \mathbf{\HiggsLambdaVal}
\end{equation}

\vspace{0.5em}
\begin{tabular}{ll}
\textbf{Prediction:} & $\HiggsLambdaVal$ \\
\textbf{Experiment:} & \HiggsLambdaExperimentalValue \\
\textbf{Accuracy:}   & $\HiggsLambdaSigma \sigma$ \\
\end{tabular}
\vspace{0.5em}

\subsection{Toll: Instability (\texorpdfstring{$\mu^2$}{mu2})}
In the Standard Model, the negative mass-squared parameter ($-\mu^2$) is an axiomatic input required to drive Spontaneous Symmetry Breaking. In the $E_8$-Persistence theory, this corresponds to the \textbf{Toll}: the entropic cost required to maintain the unstable, unbroken symmetric state ($\phi=0$) against the vacuum.

\subsubsection{The Instability Action}
As established by the topological boundary ($\chi=2$), the Electroweak interface operates as an $SU(2)$ doublet, requiring two independent complex degrees of freedom. Under \textbf{Rule 1: Capacity Partitioning}, each of these degrees of freedom must pay the self-regulation cost ($\lambda$) derived in the previous section. The total instability action ($T_H$) mandated by the substrate is therefore:
\begin{equation}
T_H = \chi \cdot \lambda = 2\lambda \approx 0.259
\end{equation}

\subsubsection{Structural Consistency: Recovering \texorpdfstring{$\mu^2 = \lambda v^2$}{mu2 lambda v2}}
We validate this entropic cost by calculating the physical mass parameter $\mu^2$. $\mu^2$ is the fractional Toll ($T_H$) scaled across the available capacity ($v$) of the interface. Dividing the total action by the two degrees of freedom yields:
\begin{equation}
\mu^2 = \frac{T_H}{2} v^2 = \frac{2\lambda}{2} v^2 = \lambda v^2
\end{equation}
This acts as a rigorous \textbf{Consistency Check}. By recovering the standard physical relation ($\mu^2 = \lambda v^2$) purely from geometric invariants, we demonstrate that the ``tachyonic'' nature of the Higgs potential is the literal manifestation of the entropic Toll driving the system toward a stable Quiescent Equilibrium.

\subsection{The Resolution Floor: The Electron Yukawa (\texorpdfstring{$y_e$}{ye})}
\label{subsec:margin_electron}
Under \textbf{Rule 2: Reversible Encoding}, any information that cannot be perfectly discretized onto the lattice grid must be stored as structural metadata to avoid dissipative loss (Landauer erasure). This metadata manifests physically as mass. The \textbf{Electron Yukawa} ($y_e$) represents the \textbf{Margin}: the absolute resolution floor of the vacuum, below which no distinct topological state can be resolved.

\subsubsection{Derivation: The Diluted Bit}
On a discrete substrate, the smallest possible signal is a single unit of action (one bit). To find the entropic cost of this signal, we calculate the dilution of this "Unit Bit" across the total configuration space of a topological defect. As established in \cite{meyer_informational_2026}, this volume is defined by the product of the defect's address budget ($L_{embed}$), its interaction aperture ($\sigma+1$), and its causal resonance area ($\Delta^2$):

\begin{equation}
y_{e,bare} \equiv Z_{MAR} = \frac{1}{V_{config}} = \frac{1}{\underbrace{L_{embed}}_{\text{Budget}} \cdot \underbrace{(\sigma+1)}_{\text{Aperture}} \cdot \underbrace{\Delta^2}_{\text{Area}}}
\end{equation}

\subsubsection{Physical Interpretation: The Landauer Limit}
Physically, the electron represents the \textbf{Landauer Limit} of the vacuum. Any topological knot with a coupling $y < y_{e,bare}$ would possess a binding energy lower than the substrate's thermal noise floor ($E < k_B T_{eff} \ln 2$). Such a state would be informationally indistinguishable from vacuum fluctuations and would immediately evaporate. The electron is stable because it is the \textbf{Minimum Viable Particle}, the physical instantiation of the smallest non-erasable unit of mass.

\subsubsection{Geometric Renormalization: The Projection Stress}
Standard QFT treats the difference between ``bare'' and ``physical'' mass as a result of virtual particle loops. In the $E_8$-Persistence framework, this correction is the \textbf{Geometric Stress} of embedding a high-dimensional knot into a lower-dimensional manifold.

A topological defect possesses an internal interaction structure of rank $\sigma=5$, but it must persist within a spacetime manifold of rank $D=4$. This mismatch ($\sigma > D$) creates an irreducible ``Overpressure'' on the surrounding lattice. We define the \textbf{Projection Coefficient} ($\mathcal{P}$) as the ratio of this internal load to the manifold's capacity:
\begin{equation}
\mathcal{P} = \frac{\text{Interaction Rank}}{\text{Manifold Rank}} = \frac{\sigma}{D} = 1.25
\end{equation}

The physical effective coupling ($y_{e,phys}$) is the bare floor enhanced by this projection stress acting through the vacuum impedance ($\alpha$):
\begin{equation}
y_{e,phys} = y_{e,bare} \cdot \left(1 + \mathcal{P} \cdot \alpha \right) = y_{e,bare} \cdot \left(1 + 1.25\alpha \right)
\end{equation}
Consistent with the dimensional budget of the Entropic Action \cite{meyer_ieaction_2026}, higher-order feedback terms are suppressed by $\alpha^2 \approx 5 \times 10^{-5}$ and are safely truncated.

\subsubsection{The Physical Result}
Substituting the invariants ($L_{embed}=31, \sigma+1=6, \Delta=43$) and the derived $\alpha$:
\begin{equation}
\begin{split}
y_{e,phys} &\approx \left(\frac{1}{31 \cdot 6 \cdot 43^2}\right) \cdot \left(1 + \frac{1.25}{137.036}\right) \\
&\approx \mathbf{\ElectronYukawaVal}
\end{split}
\end{equation}

\vspace{0.5em}
\begin{tabular}{ll}
\textbf{Prediction:} & $\ElectronYukawaVal$ \\
\textbf{Experiment:} & \ElectronYukawaExperimentalValue \\
\textbf{Accuracy:}   & $\ElectronYukawaSigma \sigma$ \\
\end{tabular}
\vspace{0.5em}

\subsubsection{Physical Consequences: The Spectral Floor}
This geometric derivation explains two fundamental features of the particle spectrum:
\begin{itemize}
    \item \textbf{Electron Stability:} The electron sits exactly at the minimum resolvable mass. There is no lower energy state for it to decay into while still maintaining its topological identity.
    \item \textbf{Neutrino Mass:} Massive neutrinos possess Yukawa couplings $y_\nu \ll y_e$. This pushes them below the vacuum's resolution floor, meaning they cannot acquire mass via the standard Higgs mechanism (which terminates at $y_e$). Their mass must arise from a secondary, non-topological mechanism (derived in future work).
\end{itemize}

% The Resonance of the Boundary. If the VEV is the "wall" containing the data, the Higgs Mass is the sound that wall makes when you strike it. It is the geometric stiffness of the capacity bound.

\subsection{Output: The Regulatory Excitation (The Higgs Mass, \texorpdfstring{$m_H$}{mH})}

Standard model-building treats the Higgs boson as a fundamental particle akin to quarks or leptons, albeit with scalar properties. In the $E_8$-Persistence framework, the Higgs is fundamentally distinct: it is not a Topological Knot (matter), but rather \textbf{Geometric Infrastructure}. It is the physical manifestation of the Resonant Interface itself.

If the Vacuum Expectation Value ($v$) represents the total state-space \textit{capacity} allocated to the interface, and the self-coupling ($\lambda$) represents the \textit{stiffness} of that partition, then the Higgs Mass ($m_H$) is the energetic cost of perturbing this regulatory infrastructure. It is the resonant vibration of the boundary.

Because the Higgs is not a free particle but the structural response of the vacuum, its mass is strictly dictated by the closure of the geometric system. By combining the interface capacity ($v_{phys}$) and the regulatory protocol ($\lambda$) derived in the preceding sections, the physical mass emerges without requiring a separate geometric derivation:

\begin{equation}
\begin{split}
m_H &= \sqrt{2\lambda} \cdot v_{phys} \\
&= \sqrt{2 (0.129424)} \cdot (246.2196 \text{ GeV}) \\
& \approx \mathbf{\HiggsMassVal}
\end{split}
\end{equation}

\vspace{0.5em}
\begin{tabular}{ll}
\textbf{Prediction:} & $\HiggsMassVal$ \\
\textbf{Experiment:} & \HiggsMassExperimentalValue \\
\textbf{Accuracy:}   & $\HiggsMassSigma \sigma$ \\
\end{tabular}
\vspace{0.5em}

\textbf{Architectural Conclusion:} This formulation resolves the ontological status of the Higgs boson. It does not exist alongside fermions in a spectral scan of lattice excitations because it is the "stage" rather than the "actor." Its mass is not an arbitrary parameter subject to infinite radiative corrections; it is the mathematically rigid impedance of the vacuum's self-regulatory partition.


\subsection{Validation: Global Impedance Matching}
The structural test of the Resonant Interface is whether it successfully satisfies \textbf{Rule 3: Impedance Matching}. For a nested subsystem to persist, its internal regulation must perfectly match the geometric aperture of the signals it governs.

We verify this by comparing the \textbf{Higgs Impedance} (the controller) to the \textbf{Weak Projected Aperture} (the load).

\begin{enumerate}
    \item \textbf{The Controller ($Z_H$):} We define the Higgs Impedance as the inverse of its regulation protocol ($1/\lambda$). Crucially, because this regulator operates under the Entropic Action filter  \cite{meyer_ieaction_2026}, its effective impedance is exponentially modulated by its action cost ($T_H = 2\lambda$):
    \begin{equation}
    Z_H = \frac{1}{\lambda} e^{-2\lambda} \approx \frac{1}{0.129424} (0.7719) \approx \mathbf{\HiggsImpedanceVal}
    \end{equation}
    
    \item \textbf{The Aperture ($A_{weak}$):} The hardware aperture of the Weak interaction is the sum of the lattice primitives involved in a state transition: the interaction links ($\sigma=5$) plus the node occupancy bit ($+1$). Because this aperture is projected onto the physical manifold, it is subject to the \textbf{Coordinate Overhead} ($\eta$):
    \begin{equation}
    A_{weak} = (\sigma + 1) \cdot \eta \approx 6 \cdot 0.994186 \approx \mathbf{\WeakApertureProjVal}
    \end{equation}
\end{enumerate}

\paragraph{The Result}
The regulator impedance ($Z_H \approx 5.9644$) matches the projected aperture ($A_{weak} \approx 5.9651$) to within \textbf{0.01\%}. 

This convergence confirms that the \textbf{Electroweak Scale} is the unique energy magnitude required to satisfy global impedance matching between the lattice resonance and the surface of matter. The $O(10^{-4})$ residual is numerically consistent with the $O(\alpha^2)$ truncation limits of the framework's geometric corrections ($\alpha^2 \approx 0.5 \times 10^{-4}$). 

The difference represents the \textbf{maximum precision} allowed by a first-order geometric derivation and identifies the theoretical floor of the vacuum's entropic drag. In the IE framework, this residual is the \textbf{minimum entropic cost} required for information to cross the interface boundary. This structural lock demonstrates that the Resonant Interface is a mathematically closed, self-sustaining regulatory subsystem, resolving the Hierarchy Problem as a necessary condition for substrate persistence.